\section{Versión con CUDA Streams}
\subsection{Motivación}
En la versión CUDA las copias D2H ocupan entre el 25 y el 30\% del tiempo y la escritura entre el 62 y el 71\%, y las dos ocurren en serie con el kernel. Surge entonces la pregunta vista en clase: ¿es posible ocultar la latencia de estas operaciones? La idea es que mientras la CPU escribe un frame, la GPU calcule y copie el siguiente por porciones, de modo que la copia de una porción se haga al mismo tiempo que el cálculo de la siguiente. Para eso CUDA provee los streams.

\begin{block}{Definición de stream}
Un stream es una cola de trabajo para el device, donde el host encola el trabajo y continúa inmediatamente.
\end{block}

\subsection{Implementación}
Sigo el esquema visto en clase (superposición kernel-D2H). Creo un vector de streams donde cada uno trabaja sobre su porción de los datos, de modo que una invocación al kernel se reemplaza por varias y cada una se superpone con su copia D2H. Divido la grilla en STREAMS = 5 bandas de filas. Lo combino con la alternativa híbrida que superpone el trabajo de la GPU con el de la CPU, que acá lo tomo como la escritura a disco.

Reservo dos grillas en el host con cudaMallocHost, porque cudaMemcpyAsync necesita memoria pinned, y así la GPU copia un frame en una mientras la CPU escribe la otra a disco.

El kernel recibe un puntero al comienzo de su banda y las filas donde empieza y termina. La fila global la sigue necesitando para detectar la frontera y las fuentes, y solo escribe las filas de su banda. El halo se lee igual que antes, porque las filas vecinas están en la misma reserva, antes o después del puntero.
\begin{lstlisting}[language=c, firstnumber=111]
__global__ void difusion(float *d_grilla, float *d_nueva, int N, int distancia, int fila_inicio, int fila_fin) {
    ...(*@\setcounter{lstnumber}{120}@*)
    int fila_banda = blockIdx.y * alto_bloque + fila_bloque;
    int fila = fila_inicio + fila_banda;
    int idx = fila_banda * N + col;
    ...(*@\setcounter{lstnumber}{145}@*)
    if (!frontera && fila < fila_fin) {
\end{lstlisting}

En cada paso, cada stream encola el kernel y la copia de su banda del frame siguiente, y mientras la GPU trabaja el host escribe a disco el frame actual. Recién al final espera a la GPU.
\begin{lstlisting}[language=c, firstnumber=214]
for (int paso = 0; paso < pasos; paso++) {
    ...(*@\setcounter{lstnumber}{218}@*)
    if (hay_siguiente) {
        for (int i = 0; i < STREAMS; i++) {
            int fila_inicio = i * N / STREAMS;
            int fila_fin = (i + 1) * N / STREAMS;
            size_t offset = (size_t)fila_inicio * N;
            ...(*@\setcounter{lstnumber}{227}@*)
            difusion<<<grid_banda, bloque, bytes_compartidos, streams[i]>>>(&d_grilla[actual][offset], &d_grilla[siguiente][offset], N, distancia, fila_inicio, fila_fin);(*@\setcounter{lstnumber}{231}@*)
            ...
            HANDLE_ERROR(cudaMemcpyAsync(&h_grilla[siguiente][offset], &d_grilla[siguiente][offset], bytes_banda, cudaMemcpyDeviceToHost, streams[i]));(*@\setcounter{lstnumber}{234}@*)
            ...
        }
    }
    ...(*@\setcounter{lstnumber}{239}@*)
    guardar_paso(h_grilla[actual], paso, N);
    ...(*@\setcounter{lstnumber}{243}@*)
    HANDLE_ERROR(cudaDeviceSynchronize());
    ...(*@\setcounter{lstnumber}{254}@*)
}
\end{lstlisting}

Ningún buffer se escribe mientras otra operación lo usa: en un paso todas las bandas leen la misma grilla y escriben filas distintas de la otra, y la copia va a la grilla del host que la CPU no está escribiendo. Como por el halo cada banda lee filas que calculan otros streams, el cudaDeviceSynchronize del final asegura que el paso anterior terminó completo antes de lanzar el siguiente.

\subsection{Resultados}
Verifiqué que los frames sean idénticos byte a byte a los de la versión sin streams, para N = 512, 1024 y 2048.

\begin{lstlisting}[language={}, numbers=none]
== Version sin streams (referencia): N=1024 pasos=30 hilos=256
== Version con streams: difusion_calor_cuda_streams.cu
...
Frames identicos (30 archivos)
\end{lstlisting}

Volví a medir el secuencial y la versión sin streams en la misma sesión, con 100 pasos, 3 repeticiones y los mismos hilos por bloque, por eso los tiempos no coinciden exactamente con los de antes. Además mido cuánto espera el host a la GPU.

\begin{table}[H]
  \centering
  \begin{tabular}{rrrrrrr}
    \toprule
    N & Hilos/bloque & Sin streams (s) & Con streams (s) & Mejora & Escritura (s) & Espera (s) \\
    \midrule
    512  & 256 & 0.0285 & 0.0209 & 1.36 & 0.0180 & 0.0001 \\
    1024 & 128 & 0.0967 & 0.0719 & 1.35 & 0.0688 & 0.0001 \\
    2048 & 128 & 0.3629 & 0.2553 & 1.42 & 0.2517 & 0.0001 \\
    4096 & 64  & 1.4108 & 0.9722 & 1.45 & 0.9657 & 0.0002 \\
    \bottomrule
  \end{tabular}
  \caption{Tiempo total de las dos versiones CUDA (100 pasos, promedio de 3 repeticiones). La escritura y la espera son de la versión con streams.}
  \label{tab:streams}
\end{table}

El tiempo total baja entre un 26 y un 31\% y la espera es prácticamente nula, o sea que la GPU queda oculta detrás de la escritura. Desde N = 1024 el tiempo total queda a menos de un 5\% del de escritura. Con N = 512 queda un 16\% por encima, porque manejar 5 streams por paso (encolar kernels y copias y leer sus eventos) tiene un costo fijo en el host que con grillas chicas ya no es despreciable.

\begin{table}[H]
  \centering
  \begin{tabular}{rrrrr}
    \toprule
    N & Secuencial (s) & Speedup sin streams & Speedup con streams & Cota \\
    \midrule
    512  & 0.0222 & 0.78 & 1.06 & 1.23 \\
    1024 & 0.0719 & 0.74 & 1.00 & 1.05 \\
    2048 & 0.3271 & 0.90 & 1.28 & 1.30 \\
    4096 & 1.9719 & 1.40 & 2.03 & 2.04 \\
    \bottomrule
  \end{tabular}
  \caption{Speedup del tiempo total frente al secuencial. La cota es el tiempo del secuencial dividido el tiempo de escritura con streams.}
  \label{tab:speedup_streams}
\end{table}

Desde N = 2048 el speedup queda prácticamente en la cota que impone la escritura, y con N = 512 se aleja por el costo de manejar los streams. Con N = 1024 empata con el secuencial porque la versión CUDA escribe los frames más lento, y la escritura sola ya tarda casi lo mismo que todo el secuencial. Por lo mismo, optimizar la GPU ya no mejora el tiempo total: con N = 4096 da entre 0.972 y 0.978 s para todas las configuraciones de hilos, incluso con 1024, cuyo kernel es 1.76 veces más lento.

La mejora frente a la versión sin streams viene de superponer la escritura a disco con el trabajo de la GPU. Una versión anterior con solo dos streams (uno de cálculo y otro de copia) ya lograba lo mismo. Dividir la grilla en bandas no agrega ganancia en este problema, como se analiza en la siguiente sección de uso de IA.